% ================================================================== COVER
\begin{titlepage}
\thispagestyle{empty}
\begin{tikzpicture}[remember picture,overlay]
  \fill[navy] (current page.north west) rectangle ([yshift=-10.2cm]current page.north east);
  \fill[teal] ([yshift=-10.2cm]current page.north west) rectangle ([yshift=-10.55cm]current page.north east);
  \fill[orange] ([yshift=-10.55cm]current page.north west) rectangle ([xshift=5.2cm,yshift=-10.75cm]current page.north west);
  % network motif
  \begin{scope}[shift={([xshift=-5.2cm,yshift=-4.6cm]current page.north east)},every node/.style={circle,fill=white!70!navy,inner sep=1.6pt}]
    \node (a) at (0,2.2) {}; \node (b) at (1.4,1.6) {}; \node (c) at (2.6,2.9) {}; \node (d) at (-1.2,0.9) {};
    \node (e) at (0.4,0.2) {}; \node (f) at (2.0,0.5) {}; \node (g) at (3.3,1.3) {}; \node (h) at (1.2,-1.0) {};
    \node (i) at (2.7,-0.9) {}; \node (j) at (3.6,3.5) {}; \node (k) at (-0.3,-0.9) {};
    \foreach \u/\v in {a/b,b/c,b/e,d/e,e/f,f/g,c/g,f/h,h/i,g/i,c/j,d/a,e/k,k/h}
      \draw[white!55!navy,line width=0.6pt] (\u)--(\v);
  \end{scope}
  \node[anchor=north west,text width=12cm,text=white,font=\sffamily] at ([xshift=2.4cm,yshift=-2.2cm]current page.north west)
   {{\bfseries\large YAOUNDÉ INSTITUTE OF BUSINESS STUDIES}\\[1pt]{\small (YIBS)}\\[14pt]
    {\fontsize{34}{38}\selectfont\bfseries Business Process\\[2pt] Analytics}\\[12pt]
    {\large Forecasts, Inventory, Optimisation and\\ Decisions under Uncertainty}};
\end{tikzpicture}
\vspace*{10.6cm}

\noindent{\sffamily
{\color{navy}\bfseries\large COMPLETE LECTURE NOTES \,--\, SECOND EDITION}\\[4pt]
{\small Seven-week practical course (7 $\times$ 4 hours) \quad$\bullet$\quad 28 contact hours \quad$\bullet$\quad capstone defence in examination week}\\[16pt]
\renewcommand{\arraystretch}{1.25}
\begin{tabular}{@{}>{\color{teal}\bfseries}l@{\hspace{1.2cm}}l@{}}
Programme & Master of Technology 2 (MTech 2)\\
Specialisation & Data Science and Artificial Intelligence\\
Institution & Yaoundé Institute of Business Studies (YIBS)\\[6pt]
Lecturer & \textbf{FOLEFAC Martins}\\
Telephone & 673\,909\,069\\
Edition & Second edition, November 2026
\end{tabular}}

\vfill
\begin{plainnote}
\textbf{Learning by doing, on real data.} Every week pairs the Mboa Fresh Market teaching case (synthetic, with
planted truths you can check) with a real public dataset: 1.07 million retail transactions, a municipal permit
process log, an industrial machine that failed, a 64\,000-customer randomised experiment and a policy change
across 27 US states. Every line of code in these notes was run end to end before printing; the output you see
is what it produced. All software is free and runs in Google Colab or on a laptop.
\end{plainnote}
\end{titlepage}

% ================================================================== CONTENTS
\tableofcontents
\clearpage
\pagenumbering{arabic}

% ================================================================== WHAT CHANGED
\chapter*{What Is New in the Second Edition}
\addcontentsline{toc}{chapter}{What Is New in the Second Edition}
\markboth{What Is New in the Second Edition}{}

The first edition was reviewed against what a strong process-analytics course should achieve for MTech Data
Science students who already know Python and machine learning. The review found a good topic list held back by
sequence, depth and data. This edition is built around one idea: \textbf{a forecast is only worth what the decision
built on it is worth}. Each week now hands its output to the next, from clean data to a forecast, from the forecast
to a stock policy, from the policy to an optimised plan, and from the plan to a test of whether it worked.

\begin{center}\small
\begin{tabularx}{\linewidth}{@{}L{4.6cm}Y@{}}
\toprule
\textbf{Finding from the review} & \textbf{What this edition does}\\\midrule
Forecasting (old Week 3) was separated from the inventory decisions that use it (old Week 6). &
Inventory moves to Week 4, directly after forecasting. Safety stock is computed from the forecast errors measured
in Lab 3, and the newsvendor order is set from a forecast quantile. \emph{Predict-then-optimise} is the spine of
the course, not a subsection.\\
Week 1 spent time on generic framing and pandas review. &
Generic material is condensed into one page and a pre-course readiness check. Week 1 goes straight to KPIs,
data quality and SQL on 1.07 million real transactions.\\
Chapters 2, 4 and 7 were grab-bags; Power BI needed Windows and a licence. &
SQL now serves the analysis instead of being a separate tool tutorial. Looker Studio (free, browser-based)
replaces Power BI as a take-home. Process mining, control charts and anomaly detection are unified as
\emph{process diagnostics}. Predictive maintenance becomes a capstone option. The capstone has its own chapter.\\
Optimisation, simulation and causal inference were thin. &
Optimisation gets a full week with a national distribution-network design. Queueing, discrete-event simulation
and stochastic optimisation share Week 6. Causal evaluation (A/B, adjustment, doubly robust, DiD with event study)
fills Week 7.\\
The capstone started in the last week; the defence was worth 10\%. &
The capstone launches in Week 2 with checkpoints in Weeks 3 and 5. The defence is worth 20\% and is held in
examination week. A learning-outcome-to-assessment map is given.\\
Only synthetic data; the generator revealed the planted patterns. &
Every week now has a real public dataset with documented licence and messiness. The Mboa generator stays, but
only to check whether a method recovers a known truth.\\
Missing classics: SPC, bullwhip, optimisation under uncertainty, SLAs. &
All four are added: $p$-charts and Phase I/II control (Week 2), the bullwhip effect and a beer-game simulation
(Week 4), two-stage stochastic programming with VSS and EVPI (Week 6), and SLA-based staffing (Week 6).\\
Overlap with machine-learning courses. &
Model building is kept brief. The new material is the decision layer: cost asymmetry, constraints, uncertainty,
service levels and causal evaluation. LightGBM appears once, as a benchmark that must beat simple baselines.\\
``The AI Operations Analyst'' risked reading as hype. &
Replaced by a concrete section on what an LLM can and cannot do in a replenishment workflow, with documented
failure cases and a tested guardrail that checks an AI-proposed order before a human approves it.\\
\bottomrule
\end{tabularx}
\end{center}

\paragraph{Tools} Everything is free and needs minimal setup: Google Colab (or local Jupyter with \texttt{uv}),
GitHub, DuckDB, statsforecast, PuLP with CBC, OR-Tools, SimPy, PM4Py and Looker Studio. Versions are pinned in
\texttt{requirements.txt}; two real version conflicts found while testing these notes are documented where they
bite (Weeks 1 and 5).
